\section{Proofs}

\subsection{Preliminaries}

\begin{definition}[Submodular Maximization~\cite{lovasz1983submodular}]\label{thm:sub_def}
	Let $\Omega$ denote a finite set, and let $f: 2^{\Omega} \rightarrow \mathbb{R}_{\geq 0}$ be a set function, where $2^{\Omega}$ is the power set of $\Omega$. The function $f$ is called submodular if it satisfies one of the following three equivalent conditions:
	\begin{itemize}
		\item For every $X, Y \subseteq \Omega$ with $X \subseteq Y$, and for all $x \in \Omega \setminus Y$, we have 
			\begin{equation}
				f(X \cup \{x\}) - f(X) \geq f(Y \cup \{x\}) - f(Y) \,,
			\end{equation}
		\item For every $X, Y \subseteq \Omega$, we have 
			\begin{equation}\label{thm:sub_def_cond2}
				f(X) + f(Y) \geq f(X \cup Y) + f(X \cap Y) \,,
			\end{equation}
		\item For every $X \subseteq \Omega$ and $x_1, x_2 \in \Omega \setminus X$ such that $x_1 \neq x_2$, we have 
			\begin{equation}
				f(X \cup \{x_1\}) + f(X \cup \{x_2\}) \geq f(X \cup \{x_1, x_2\}) + f(X) \,.
			\end{equation}
	\end{itemize}
\end{definition}
These three conditions are equivalent, and the first condition is the most commonly used, as it directly reflects the law of diminishing marginal utility as the number of items increases.

\begin{definition}[Sub-additive]
	A set function $f: 2^{\Omega} \rightarrow \mathbb{R}_{\geq 0}$ is sub-additive if for every two sets $X, Y \in \Omega$, we have \
	\begin{equation}
		f(X\cup Y)\leq f(X) + f(Y)\,.
	\end{equation}
\end{definition}

\begin{lemma}\label{thm:sub_additive}
	A non-negative submodular set function $f: 2^{\Omega}\rightarrow\mathbb{R}_{\geq 0}$ is sub-additive.
	\begin{proof}
		As the second condition \cref{thm:sub_def_cond2} in \cref{thm:sub_def}. 
		for $X,Y\subseteq \Omega$, we have $f(X) + f(Y) \geq f(X \cup Y) + f(X \cap Y)$. So, $f(X) + f(Y) \geq f(X \cup Y)$ as $f(X \cap Y) \geq 0$.
	\end{proof}
\end{lemma}

\begin{lemma}\label{thm:fs_submodular}
	Let $f: 2^{\Omega}\rightarrow \mathbb{R}$ be submodular. 
	Let $S\subseteq \Omega$, and $f_{S}(X)=f(S\cup X)-f(S)$ for every $X\subseteq\Omega$. ($f_{S}$ is the marginal value function for set $S$.) 
	Then $f_S$ is also submodular.
	\begin{proof}
		Let $X, Y\subseteq\Omega \setminus S$; it suffices to consider ground set $\Omega\setminus S$.
		\begin{align}
			&\left(f_S(X\cup Y)+f_S(X\cap Y)\right)-\left(f_S(X)-f_S(Y)\right) \notag\\
			&\quad=f(S\cup X\cup Y)-f(S)+f(S\cup(X\cap Y))-f(S) \notag\\
			&\qquad-\left(f(S\cup X)-f(S)+f(S\cup Y)-f(S)\right) \\
			&\quad = f(S\cup X\cup Y)+f(S\cup(X\cap Y)) \notag\\
			&\qquad-f(S\cup X)-f(S\cup Y) \notag\\
			&\quad \leq 0 \,. \notag
		\end{align}
		The last inequality is by $S\cup X\cup Y=(S\cup X)\cup(S\cup Y)$, $S\cup(X\cap Y)=(S\cup X)\cap(S\cup Y)$ and submodularity of $f$.
		Therefore, $f_S$ is also submodular is proved.
	\end{proof}
\end{lemma}

\subsection{Proof of $(1-1/e)$-approximation}

Submodular maximization is NP-hard in general.
Therefore, most research in this field focuses on approximation algorithms with polynomial-time complexity. 
While the submodular function is monotone, \ie, for every $X, Y \subseteq \Omega$, we have $f(X) \leq f(Y)$. 
The problem of maximizing a monotone submodular function subject to a cardinality constraint admits a $(1 - 1/e)$-approximation greedy algorithm~(as introduced  in \cref{subsec:method_dpp})~\cite{nemhauser1978analysis}.

In this section, we provide a concise proof of the $1 - 1/e$ approximation ratio for the greedy algorithm.
\begin{theorem}[$(1 - \frac{1}{e})$-approximation of Greedy Algo.]
    There exists a greedy algorithm for the submodular maximization problem, which starts with an empty set \( S = \emptyset \) and iteratively selects the item that maximizes the marginal gain:
    \begin{equation}\label{eq:greed_sub}
        j = \mathop{\mathrm{argmax}}\limits_{i \in \Omega \setminus S} \quad f(S \cup \{i\}) - f(S) \,.
    \end{equation}
    The algorithm continues until the selected set \( S \) reaches the cardinality limit \( k \).
    
    This greedy algorithm provides a solution \( \hat{S} \), which guarantees a $(1 - 1/e)$ approximation, where the optimal solution is denoted as \( S^* \):
    \begin{equation}
        f(\hat{S}) \leq f(S^*)
    \end{equation}
\end{theorem}
\begin{proof}
    According to \cref{thm:fs_submodular}, \(f_S\) is submodular, and by \cref{thm:sub_additive}, it is also sub-additive. Therefore, we have:
    \begin{equation}
        f_S(S^*) \leq \sum_{x \in S^*} f_S(x)\,,
    \end{equation}
    which implies that:
    \begin{equation}
        \exists x \in S^*, \quad f_S(x) \geq \frac{1}{k} f_S(S^*) \,.
    \end{equation}
    For this \(x\), we have the following margin lower bound:
    \begin{equation}\label{eq:margin_lowerbound}
        f(S \cup \{x\}) - f(S) \geq \frac{f(S^*) - f(S)}{k}\,.
    \end{equation}
    
    Let \( S_t \) denote the selected subset after the \(t\)-th step of the greedy algorithm. According to \cref{eq:greed_sub}, in the greedy algorithm, we have:
    \begin{equation}
        f(S_{t+1}) - f(S_t) \geq f(S_t \cup \{x\}) - f(S_t), \quad \forall x \in \Omega \setminus S_t\,.
    \end{equation}
    Therefore, in the greedy algorithm, the marginal gain is lower-bounded by:
    \begin{equation}\begin{aligned}\label{eq:margin_gain_lowerbound}
        f(S_{t+1}) - f(S_t) &\geq \frac{f(S^*) - f(S_t)}{k}\,.
    \end{aligned}\end{equation}
    This implies:
    \begin{equation}
        f(S^*) - f(S_{t+1}) \leq \left(1 - \frac{1}{k}\right)(f(S^*) - f(S_t))\,.
    \end{equation}
    Hence, when the greedy algorithm selects a subset \( \hat{S} = S_k \) after \(k\) steps, we have:
    \begin{equation}\begin{aligned}
        f(S^*) - f(S_k) &\leq \left(1 - \frac{1}{k}\right)(f(S^*) - f(S_{k-1})) \\
        &\leq \left(1 - \frac{1}{k}\right)^2(f(S^*) - f(S_{k-2})) \\
        & \qquad \vdots \\
        &\leq \left(1 - \frac{1}{k}\right)^k(f(S^*) - f(S_0))\,,
    \end{aligned}\end{equation}
    where \( S_0 \) is the initial set at \( t = 0 \), with \( S_0 = \emptyset \), such that \( f(S_0) = 0 \).
    Therefore, we obtain:
    \begin{equation}
        f(S^*) - f(\hat{S}) \leq \left(1 - \frac{1}{k}\right)^k f(S^*)\,.
    \end{equation}
    Hence, we have:
    \begin{equation}\begin{aligned}\label{eq:detailed_bound}
        f(\hat{S}) 
            &\geq \left( 1 - \left(1 - \frac{1}{k}\right)^k \right) f(S^*) \\
            &\geq \lim_{k \to +\infty} \left( 1 - \left(1 - \frac{1}{k}\right)^k \right) f(S^*) \\
            & = \left(1 - \frac{1}{e}\right) f(S^*)\,.
    \end{aligned}\end{equation}
    The proof is complete.
\end{proof}


Previous works~\cite{fujishige2005submodular,krause2014submodular} have established that the determinantal point process~(DPP) is a monotone submodular function.
Therefore, when selecting 8 frames using the standard DPP, the approximation ratio is at least \mbox{\( 1 - \left( 1 - \frac{1}{8} \right)^8 = 65.6\% \).}


\subsection{Proof of Time Complexity}
The time complexity of \cref{alg} is closely $\mathcal{O}(nk^3)$, where $n$ represents the number of candidate frames and $k$ denotes the selection size.


\begin{proof}
    As indicated by the pseudocode in \cref{alg}, there are three loops, with the iteration sizes specified as follows:
    \begin{itemize}
        \item Line 3: $\lceil \frac{n}{m} \rceil$;
        \item Line 6: $k + 1$;
        \item Line 7: $k_{t-1} \leq \min(m, k, k - C_{t-1}) \leq \min(m, k)$.
    \end{itemize}
    Here, $m$ denotes the segment size. 
    The DPP update in each iteration (Line 9) has a time complexity of $\mathcal{O}(mk_{t-1})$, utilizing Cholesky decomposition for incremental computation.

    Therefore, the overall time complexity can thus be expressed as:
    \begin{equation}
        \mathcal{O}\left(
            \frac{n}{m} \cdot (k \cdot mk_{t-1}^2 ) 
        \right) = \mathcal{O}\left(
            n \cdot \min(m^2k, k^3)  
        \right)
    \end{equation}
    Additionally, the time complexity of computing the determinant in line 5 is $\mathcal{O}(m^3)$, resulting in the following time complexity:
    \begin{equation}
        \mathcal{O}\left(
            \frac{n}{m} \cdot m^3
        \right) = \mathcal{O}\left(
            n m^2
        \right)
    \end{equation}
    Next, we analyze the relationship between $m$ and $k$, leading to the following total time complexity:
    \begin{equation}
        \left\{
            \begin{array}{ll}
                \mathcal{O}(nm^2),  & \text{if } k < m^{2/3} \\
                \mathcal{O}(nk^3),  & \text{if } m^{2/3} \leq k < m  \\
                \mathcal{O}(nm^2k) < \mathcal{O}(nk^3), & \text{if } k \geq m \\
            \end{array}
        \right.
    \end{equation}
    In practice, the segment size $m$ is typically small, so in most practical cases, the time complexity of \cref{alg} is upper-bounded by $\mathcal{O}(nk^3)$.
\end{proof}
This algorithm exhibits pseudo-polynomial time complexity, analogous to the knapsack problem. 
In the dynamic programming approach proposed in this paper, the runtime of the pseudo-polynomial complexity is practically comparable to that of a polynomial-time algorithm.

In practical applications, the segment size $m$ and selection size $k$ are both much smaller than the total number of video frames $n$, \ie, $m \ll n$ and $k \ll n$. 
Thus, regardless of the relationship between $m$ and $k$, the total time complexity is much smaller than feeding all frames into transformer-based LLMs, which have a time complexity of $O((n \cdot \text{\#tokens\_per\_image})^2)$ per layer and attention head.

Additionally, the iteration in line 6 of \cref{alg} is independent, enabling parallel updates. 
This results in a time complexity of \(\mathcal{O}(nk^2)\), making its efficiency comparable to that of the standard DPP. 
Moreover, without lazy strategies and parallel updates, the time complexity is closely \(\mathcal{O}(nk^4)\) in most practical cases, and the proof is similar.


